\FloatBarrier
\section*{Выводы по главе 2}
\phantomsection\addcontentsline{toc}{section}{Выводы по главе 2}

В главе~2 сформулирована единая математическая основа для численного
моделирования смазочного слоя в подшипниках скольжения как для гладкой,
так и для детерминированно микроструктурированной рабочей поверхности.

\begin{enumerate}[wide, labelindent=\parindent, nosep]
  \item Для трёх типовых постановок (упорной, радиальной
    и возвратно-поступательной) введены согласованные поверхностные
    координаты, масштабы безразмеризации и унифицированная форма
    уравнения Рейнольдса, что обеспечивает сопоставимость результатов
    и возможность переноса алгоритмов между постановками.
  \item Сформулирована параметрическая модель детерминированного
    микрорельефа в виде добавки к базовой функции зазора:
    $h = h_{\mathrm{base}} + \Delta h$, где $\Delta h$ задаётся суммой
    вкладов локализованных элементов.  Введены основные геометрические
    параметры элемента (глубина~$h_p$, полуоси~$a$, $b$,
    ориентация~$\alpha$, параметр плоского дна~$\rho_0$,
    длина канавок~$\ell$).
  \item Систематизированы типы геометрии микроструктур
    (эллипсоидальная, параболическая, коническая, косинусная, smoothstep,
    с плоским дном -- для эллиптических элементов; поперечный
    и продольный профили -- для канавок).
    Описаны типы раскладки элементов на поверхности (прямоугольная,
    шахматная, гексагональная, спиральная по принципу филлотаксиса,
    а также раскладка канавок) и сформулировано правило включения
    элементов в область текстурирования.
  \item Введён компактный набор безразмерных параметров микрорельефа
    ($H_p$, $A^{*}$, $B^{*}$, $\kappa$, $P_1^{*}$, $P_2^{*}$, $\phi$,
    $\rho_0$, $\ell^{*}$), позволяющий проводить параметрические
    исследования и сравнивать различные конфигурации при одинаковых
    безразмерных параметрах.
  \item Рассмотрена физическая природа кавитации и показаны ограничения
    инженерного подхода отсечения давления снизу (условие Гюмбеля),
    связанные с нарушением баланса массы.  Для корректного учёта разрыва
    плёнки сформулирована массосохраняющая постановка кавитации (JFO),
    включающая уравнение
    сохранения для~$\vartheta h$ и условия комплементарности
    для~$p$ и~$\vartheta$.
  \item Описана связь постановки JFO с численной реализацией:
    дискретизация в консервативной форме и итерационное решение методом
    активного множества.  Подробности дискретизации потоков, выбора
    решателя и проверки массобаланса вынесены в главу~3.
\end{enumerate}

Сформулированные модели, параметризация микрорельефа и массосохраняющая
модель кавитации (JFO) образуют основу для реализации расчётного алгоритма
и проведения вычислительных экспериментов, представленных в последующих
главах.
